% $Header: /home/vedranm/bitbucket/beamer/solutions/conference-talks/conference-ornate-20min.en.tex,v 90e850259b8b 2007/01/28 20:48:30 tantau $

\documentclass{beamer} %[compress]
% alternatively add handout or draft option

%\mode<handout>%{\setbeamercolor{background canvas}{bg=black!5}}
%\usepackage{pgfpages}
%\pgfpagesuselayout{4 on 1}[a4paper,border shrink=5mm,landscape]
% For handout printing, 4 slides on 1 page

\mode<presentation>
{
  \usetheme{Pittsburgh} % plain template Pittsburgh
  \useoutertheme{infolines} % bottom (author, title, place)
  \useoutertheme[subsection=false]{miniframes} % top (sections, frames)
  \usecolortheme{beaver} %blue color whale
	\setbeamertemplate{frametitle}[default][right] % frame titles on the right
  \setbeamercovered{transparent}
  % or whatever (possibly just delete it)
}

\setbeamersize{text margin left=1cm, text margin right=1cm}
%\setbeamertemplate{navigation symbols}{}
\usepackage[english]{babel}
\usepackage{booktabs}
\usepackage{comment}
% or whatever

\usepackage[cp1250]{inputenc}
% or whatever
\usepackage{subfigure} %subfigures

\usepackage{times}
\usepackage{tikz}
%\usepackage{breqn}
\usepackage{amsmath}
\usepackage[T1]{fontenc}
\def\sym#1{\ifmmode^{#1}\else\(^{#1}\)\fi} %shortcut for Stata tables
% Or whatever. Note that the encoding and the font should match. If T1
% does not look nice, try deleting the line with the fontenc.

%\logo{\includegraphics[height=0.5cm]{logo.jpg}} % logo on every page



\title[Lecture 7: Publication Bias] % (optional, use only with long paper titles)
{Lecture 7: Publication Bias}

\author[T.\ Havranek] % (optional, use only with lots of authors)
{ Tomas Havranek
}
% - Give the names in the same order as the appear in the paper.
% - Use the \inst{?} command only if the authors have different
%   affiliation.

\institute[Prague] % (optional, but mostly needed)
{
  	Charles University, CEPR, METRICS

  %\inst%
  %Economic Research Department\\
  %Czech National Bank
      %\inst{1}%
  %Research Department\\
  %Czech National Bank
  %\and
  %\inst{2}%
  %Institute of Economic Studies\\
  %Charles University, Prague
  
%\pgfdeclareimage[height=1cm]{logo}{logo.jpg} % logo only on title page
%\pgfuseimage{logo}
}% - Use the \inst command only if there are several affiliations.
% - Keep it simple, no one is interested in your street address.

\date[Meta in Econ and Finance] % (optional, should be abbreviation of conference name)
{Research Synthesis in Economics and Finance, \\ University of Canterbury\\[2ex]
  March 10, 2025
}%\\[4mm]% - Either use conference name or its abbreviation.
% - Not really informative to the audience, more for people (including
%   yourself) who are reading the slides online


%\subject{Intertemporal Substitution}

\begin{comment}
\AtBeginSection[]
{
  \begin{frame}<beamer>{Outline}
    \tableofcontents[currentsection]%,currentsubsection]
  \end{frame}
}
\end{comment}

\begin{document}

\begin{frame}
  \titlepage
\end{frame}

\begin{comment}

\begin{frame}{Outline}
  \tableofcontents
  % You might wish to add the option [pausesections]
\end{frame}

\end{comment}

%\section{Summary}
%\subsection{}

%\section{EIS}
%\subsection{}


\begin{frame}{Intuition}
\begin{figure}
	\centering
		\includegraphics[width=1.00\textwidth]{c1.png}
\end{figure}
\end{frame}

\begin{frame}{Intuition}
\begin{figure}
	\centering
		\includegraphics[width=0.9\textwidth]{c2.png}
\end{figure}
\end{frame}

\begin{frame}{Intuition}
\begin{figure}
	\centering
		\includegraphics[width=0.9\textwidth]{c3.png}
\end{figure}
\end{frame}

\begin{frame}{Intuition}
\begin{figure}
	\centering
		\includegraphics[width=0.51\textwidth]{c4.png}
\end{figure}
\end{frame}

\begin{frame}{Paxil scandal}
    \begin{itemize}
        \item \textbf{Background:}
        \begin{itemize}
            \item Paxil (paroxetine): Antidepressant approved for adolescents, developed by GlaxoSmithKline (GSK).
            \item Study 329 (1998-2001): Clinical trial funded by GSK to test Paxil's safety and efficacy in adolescents.
        \end{itemize}
				
				\bigskip
        \item \textbf{What Happened?}
        \begin{itemize}
            \item Published study claimed Paxil was \textit{"well-tolerated and effective."}
            \item Internal documents later revealed:
            \begin{itemize}
                \item Negative results (e.g., increased suicidal ideation) were downplayed.
                \item Positive outcomes were selectively reported.
            \end{itemize}
        \end{itemize}
    \end{itemize}
\end{frame}

\begin{frame}{Paxil scandal}
    \begin{itemize}
        \item \textbf{Discovery of Bias:}
        \begin{itemize}
            \item Legal proceedings in 2012 forced disclosure of internal data.
            \item Independent reanalysis of Study 329 showed:
            \begin{itemize}
                \item No significant benefit over placebo.
                \item Serious safety concerns (e.g., suicidal ideation).
            \end{itemize}
        \end{itemize}
				
				\bigskip
        \item \textbf{Key Lessons:}
        \begin{itemize}
            \item Publication bias can harm patients and public trust.
            \item Highlights the need for:
            \begin{itemize}
                \item Transparency in research data.
                \item Independent replication and reanalysis.
            \end{itemize}
        \end{itemize}
    \end{itemize}
\end{frame}

\begin{frame}{No bias, no correlation (?)}
Researchers estimate

   \begin{equation*}
    Earnings_j = \tikz[baseline={(E.base)}]{\node[draw,circle,inner sep=1pt] (E) {$\alert{\gamma}$};} Beauty_j + w_j.
    \end{equation*}
		
\smallskip

They assume that $\widehat{\gamma}/SE_{\widehat{\gamma}}$ has a $t$-distribution. 

\smallskip

$\rightarrow$ \alert{estimates should not be correlated with standard errors} (Card \& Krueger, 1995)

\bigskip

\begin{block}{Types of publication bias}
Type I: Selection for the ``correct sign.'' \\
Type II: Selection for statistical significance.
\end{block}

\end{frame}


\begin{frame}{Simulation}
\begin{figure}
	\centering
		\includegraphics[width=1.00\textwidth]{1.pdf}
\end{figure}
\end{frame}

\begin{frame}{Simulation}
\begin{figure}
	\centering
		\includegraphics[width=1.00\textwidth]{2.pdf}
\end{figure}
\end{frame}

\begin{frame}{Simulation}
\begin{figure}
	\centering
		\includegraphics[width=1.00\textwidth]{3.pdf}
\end{figure}
\end{frame}

\begin{frame}{Simulation}
\begin{figure}
	\centering
		\includegraphics[width=1.00\textwidth]{4.pdf}
\end{figure}
\end{frame}

\begin{frame}{Simulation}
\begin{figure}
	\centering
		\includegraphics[width=1.00\textwidth]{5.pdf}
\end{figure}
\end{frame}

\begin{frame}{Simulation}
\begin{figure}
	\centering
		\includegraphics[width=1.00\textwidth]{6.pdf}
\end{figure}
\end{frame}

\begin{frame}{Simulation}
\begin{figure}
	\centering
		\includegraphics[width=1.00\textwidth]{7.pdf}
\end{figure}
\end{frame}

\begin{frame}{Simulation}
\begin{figure}
	\centering
		\includegraphics[width=1.00\textwidth]{8.pdf}
\end{figure}
\end{frame}

\begin{frame}{Simulation}
\begin{figure}
	\centering
		\includegraphics[width=1.00\textwidth]{9.pdf}
\end{figure}
\end{frame}





\begin{frame}{Funnel asymmetry test}
In the absence of publication bias the funnel should be symmetrical:
\begin{equation*}
\mbox{estimate}_i = \underbrace{\beta}_{\mbox{true effect}} + \underbrace{\beta_0 SE_i}_{\mbox{publication bias}} + \mu_{i}.
\end{equation*}

The equation is heteroscedastic. Weighted least squares yield 

\begin{equation*} 
t_i = \beta_0 + \beta \left(\frac{1}{SE_i}\right) + \vartheta_{i}.
\end{equation*} 
This is call the funnel asymmetry test---precision effect test (FAT-PET).
\end{frame}



\begin{frame}{PEESE}
Precision effect test with standard error (PEESE):

\bigskip

\begin{equation*}
\mbox{estimate}_i = \underbrace{\beta}_{\mbox{true effect}} + \underbrace{\beta_0 SE^2_i}_{\mbox{publication bias}} + \mu_{i}.
\end{equation*}

\bigskip

The equation is heteroscedastic. Weighted least squares yield 

\begin{equation*} 
t_i = \beta_0 SE_i + \beta \left(\frac{1}{SE_i}\right) + \vartheta_{i}.
\end{equation*} 
\end{frame}



\begin{frame}{WAAP (Ioannidis et al., 2017)}

          \textbf{W}eighted \textbf{A}verage of \textbf{A}dequately \textbf{P}owered (WAAP) estimates.
					
						\bigskip

        \begin{itemize}
            \item Focus on studies with adequate statistical power to reduce the impact of selective reporting.
            \item Combines results only from studies with at least 80\% power.
        \end{itemize}
				
				\bigskip
				
       \textbf{Issues:}
        \begin{itemize}
            \item Simple and robust.
            \item Need to estimate the true effect first and then compute power.
        \end{itemize}

\end{frame}

\begin{frame}{Stem-based correction technique (Furukawa, 2019)}

\textbf{Stem-based method:}
\begin{itemize}
    \item A non-parametric estimator that exploits the variance-bias trade-off.
    \item Focuses on the most precise estimates, referred to as the ``stem'' of the funnel plot.
\end{itemize}

\bigskip

\textbf{Key features:}
\begin{itemize}
    \item Robust to various publication selection processes.
    \item Does not rely on specific distributional assumptions.
\end{itemize}

\bigskip

\textbf{Application:}
\begin{itemize}
    \item Filters studies with the highest precision (lowest standard errors) to estimate the true effect.
    \item Reduces influence from noisier, potentially biased studies in the broader funnel plot.
\end{itemize}

\end{frame}



\begin{frame}{Endogenous kink model (Bom \& Rachinger, 2019)}

\textbf{Overview:}
\begin{itemize}
    \item Fits a piecewise linear model with a kink determined by the data.
    \item Identifies where publication selection is less likely to occur.
\end{itemize}

\bigskip

\textbf{Key features:}
\begin{itemize}
    \item Endogenously estimates the cutoff based on standard errors.
    \item Often reduces bias and improves efficiency compared to traditional methods.
\end{itemize}

\bigskip

\textbf{Application:}
\begin{itemize}
    \item Highlights genuine effects by focusing on less selective segments of the data.
    \item Effective in meta-analyses with substantial heterogeneity in study precision.
\end{itemize}

\end{frame}


\begin{frame}{Selection model (Andrews \& Kasy, 2019)}

\textbf{Overview:}
\begin{itemize}
    \item Models the selection process explicitly to address publication bias.
    \item Assumes researchers are more likely to publish statistically significant results.
\end{itemize}

\bigskip

\textbf{Key features:}
\begin{itemize}
    \item Uses the likelihood of publication as a function of statistical significance.
    \item Adjusts estimates to reflect the missing studies.
\end{itemize}

\bigskip

\textbf{Comparison to funnel-based models:}
\begin{itemize}
    \item Selection models are more rigorous but require stronger assumptions about the publication process.
    \item Funnel methods are simpler and more robust to model misspecification.
\end{itemize}

\end{frame}




\begin{frame}{Example: labor supply elasticity}
\begin{figure}
	\centering
		\includegraphics[width=0.9\textwidth]{_funnel_all2.pdf}
\end{figure}
\end{frame}

\begin{frame}{Results vary: RoBMA?}
\begin{figure}
	\centering
		\includegraphics[width=1.00\textwidth]{fat.pdf}
\end{figure}
\end{frame}

\begin{frame}{Bias caused by preference for positive estimates}
\begin{figure}
	\centering
		\includegraphics[width=0.9\textwidth]{_caliper2.pdf}
\end{figure}
\end{frame}

\end{document}


